\documentclass[12pt]{article}
\usepackage[margin=1in]{geometry}
\usepackage{enumitem}
\usepackage{titlesec}
\usepackage{parskip}
\usepackage{graphicx}
\usepackage{hyperref}
\usepackage{fontenc}

\title{\textbf{Huber and Shipan 2002 RG Notes}\\
\large Huber, John D., and Charles R. Shipan. \textit{Deliberate Discretion? The Institutional Foundations of Bureaucratic Autonomy}.\\
Cambridge: Cambridge University Press, 2002.}
\author{}
\date{}

\begin{document}

\maketitle

\section{Main Idea}



%--------------------------------------------------
\section{General Notes}
%--------------------------------------------------

\begin{itemize}[leftmargin=*, itemsep=4pt]
    \item Though voters care about policy outcomes, they are generally agnostic towards the actual writing of the laws that drive them. These statutory details are core drivers of policy outcomes, but are therefore largely ignored by the general public. As a result, politicians have two general strategies in achieving policy ends through statutory means:
    \begin{enumerate}
        \item ``Write long statutes with extremely detailed language in an effort to micromanage the policymaking process''
        \item ``Write vague statutes that leave many details unspecified, thereby delegating policymaking authority to other actors, usually bureaucrats''
    \end{enumerate}
    \item ``Our argument draws attention to four general factors that affect politician's incentives and ability to use legislation to micromanage policy-making:"
    \begin{enumerate}
        \item The level of policy conflict between the politicians who adopt the statutes and the agents who implement them;
        \item The capacity of politicians to write detailed statutes;
        \item The bargaining environment in which statutes are adopted;
        \item the expectations that politicians have about nonstatutory factors (course or legislative oversight opportunities), independent of the language of the statutes
    \end{enumerate}
    \item Good governance essentially requires the construction of the \emph{administrative state}, but the administrative state itself is a somewhat undemocratic phenomenon: accountability is indirect, and power is vested in hands that are not necessarily elected or even closely impacted by elections. This is particularly notable given the high level of power/authority vested within bureaucratic systems, who are responsible for the vast majority of public policymaking in modern societies.
    \item Frame statutes as ``blueprints'' of public policy: set out the framework and guidelines that bureaucrats will subsequently follow, and in doing so choose a degree of specificity / vagueness to in turn yield some degree of autonomy.
    \item To a large extent, statutory specificity depends on \emph{legislative capacity}, which in turn ``affects the costs to politicians of drafting detailed legislation.'' Costs are information and opportunity-related (time, other resources, etc.) Highly varied at state legislature levels.
    \item ``If politicians can depend on other nonstatutory features of the political environment to help produce bureaucratic actions that are congruent with their preferences, this diminishes the extent to which precise statutory language is necessary for obtaining desired policy outcomes.''
    \item Two core mistakes possible from this arrangement:
    \begin{enumerate}
        \item \underline{Bureaucrat given too much discretion}: Bureaucrat is empowered to choose a policy yielding an outcome that matches their own ideal point, but may be far from the politician's ideal point
        \item \underline{Bureaucrat given not enough discretion}: Because bureaucrat has additional (and private) information about policy-outcome mapping, they may know that delivering the politician's ideal outcome requires enacting a policy far from the politician's ideal policy -- yet, if they are not sufficiently empowered to stray that far, they may be bounded in how close they can deliver to the ideal outcome
    \end{enumerate}
    \item Arguments as to how the level of discretion in legislation should vary with different politican and institutional contexts (all direct quotes):
    \begin{itemize}
        \item As policy conflict between politicians and bureaucrats increases, or as legislative capacity increases, the relative value to politicians of low-discretion statutes increases.
        \item Whether the model predicts a direct or interactive effect of these two variables on discretion depends on the distribution of legislative capacity and policy conflict among the political systems that are analyzed. If there are a large number of political systems that never have sufficient capacity to adopt a low-discretion law, then we should expect interactive effects (where, for example, legislative capacity influences the design of statutes only if the level of policy conflict is sufficiently high). By contrast, if a large numbere of political systems have high legislative capacity, the effects of these two variables should be direct.
        \item As the ability of politicians to rely on nonstatutory factors to achieve policy objectives increases, the relative value of low-discretion laws declines.
        \item These results hold in parliamentary and presidential systems.
        \item In presidential systems, bicameral conflict often leads to more discretion.
    \end{itemize}
    \item ``When conflict is absent and legislators trust policymakers in the agency, they are likely to design vague statutes that allow the agency to draw on its expertise to make policy. And when conflict is present, legislators who disagree with the agency's policy views will constrain the agency by putting specific wording into legislation.''
    \item ``When a different party controls each chamber, the party favoring the governor has the incentive and ability to block some -- but not all -- legislative language that ties the hands of the relative agency.''
\end{itemize}


\section{Important Specific Factual Claims}

\begin{itemize}[leftmargin=*, itemsep=4pt]
    \item ``Wise and salutary neglect'': Edmund Burke's 1775 \emph{Speech on the Conciliation of America} -- providing a broad latitude to operate independently and subsequently create individual laws, procedures, etc. Synonymous with high political autonomy
    \item A core factor in the necessity of the administrative state is an informational problem: we need technical expertise for certain tasks that cannot necessarily be centralized in the elected political system. This requires providing some large degree of autonomy, which in turn invites risk (think principal-agent problem)
    \item Model largely built on informational asymmetry: there is some mapping from policy to outcome, which is perfectly known by the bureaucrat, but is not known by the politician. Both care about outcome, but politician can therefore not perfectly match it, and requires the bureaucrat's information.
    \item Prime case study is Michigan's Medicaid expansion in the late '90s: \\
    \begin{quote}
        There was, however, an alternative to having the state set prices itself. The state realized that although it did not know the marginal costs and benefits faced by the providers, someone else did -- namely, the providers themselves. These providers knew better than anyone else what the true costs were. The question then became how to get this information from the providers. Furthermore, in order for these providers to know the costs and benefits of participation, they would have to know exactly what services they would need to provide. And this, needless to say, would require addressing all the technical issues just discussed.
    \end{quote}
    \item Another quote on single-party dominance: \\
    \begin{quotation}
        The legislator and Havemann were having a meal together in April and she said, ``We should get together and talk.'' And his reply was, ``I don't know why we need to talk to you. You have no way to stop us if we want to do something.''
    \end{quotation} 
    \item On Michigan case: ``Two features of this change in the design of statutes are worth noting. First, the constraints that appeared in the statutes were largely narrow descriptions of policy rather than procedural constraints on agency decision-making processes. Procedures were clearly present and non-trivial, but the participants we talked to emphasized the policy restrictions, which represent the lion's share of the new legislative language. Second, although the number of policy constraints rose when the Democrats took over the House, the Democrats were not able to insert all the policy constraints they desired into the statutes. Both our interviews and our examiniation of the initial House version of the MMC legislation in 1997 make it clear that several of the Democrats' desired policies had to be left out of the final version of the bill."
\end{itemize}


\section{Connections}

\begin{itemize}
    \item Where previous literature in autonomy generally focuses on a single case at a time, Huber \& Shipman provide a novel comparative perspective so that the institutional environment can be evaluated as a variable (rather than a fixed component of the situational context).
\end{itemize}


\section{My Critiques}

\begin{itemize}
    \item
\end{itemize}


\end{document}
